%% v2_methodology.tex
\chapter{V2 Agentic RAG: System Architecture}
\label{chap:v2_methodology}

This chapter presents Version 2 (V2) of the enterprise retrieval system: a fully
agentic Retrieval-Augmented Generation pipeline that extends the V1 adaptive
retrieval planner into a five-stage agent architecture with iterative
self-correction. V1 addressed retrieval strategy selection as a single-step
planning problem; V2 frames the full question-answering pipeline as an
agent loop, enabling multi-hop query decomposition, relevance-based retrieval
validation, and automatic re-retrieval when initial results are insufficient.

\section{Design Principles}
\label{sec:v2_design}

The V2 architecture is guided by three design principles derived from the
V1 failure mode analysis (Section~\ref{sec:disc_failure_modes}):

\textbf{(1) Decompose reasoning, not just routing.} V1 treated each query as
atomic, selecting one strategy and executing one retrieval call. Multi-hop
queries and aggregation tasks consistently underperformed because a single
retrieval call cannot satisfy a query that requires evidence from multiple
documents. V2 introduces a Query Analysis Agent that decomposes complex
queries into ordered sub-queries, each with its own retrieval plan.

\textbf{(2) Validate evidence, not just retrieve.} V1 returned raw ranked
results without assessing their relevance to the specific query. V2 adds a
Validation Agent that scores each retrieved chunk against its originating
sub-query text and triggers re-retrieval when the relevance signal is weak.
This transforms retrieval from an open-loop procedure into a closed-loop
feedback process.

\textbf{(3) Preserve V1 components as fallback baselines.} The V1 rule planner
remains the core of the V2 Planning Agent. LLM-based components are layered
on top and wrapped with fallback paths, so the V2 pipeline degrades gracefully
to V1 behavior when LLM providers are unavailable.


\section{Pipeline Architecture}
\label{sec:v2_pipeline}

The V2 pipeline is implemented as a LangGraph StateGraph
\citep{langgraph2024} — a directed graph of agent nodes with a shared
typed state object. Five agents are connected in sequence, with a conditional
re-retrieval back-edge after the Validation Agent:

\begin{center}
\textit{Query Analysis} $\rightarrow$ \textit{Planning} $\rightarrow$
\textit{Retrieval} $\rightarrow$ \textit{Validation}
$\xrightarrow{\text{passed}}$ \textit{Answer} $\rightarrow$ \textit{End}
\end{center}
\begin{center}
\textit{Validation} $\xrightarrow{\text{failed, loop < 3}}$
\textit{Loop Counter} $\rightarrow$ \textit{Planning} \quad (re-retrieval)
\end{center}

The shared state (\texttt{AgentState}) is a TypedDict with eleven fields.
Each agent node returns a partial dict; LangGraph merges it into the
current state. The \texttt{trace} field is annotated with
\texttt{operator.add} for append semantics, accumulating the provenance
of every agent decision. All other fields are replaced on each write, so
the \texttt{retrieval\_outputs} field always holds only the most recent
retrieval attempt. Table~\ref{tab:v2_state} summarises the key state fields.

\begin{table}[H]
\centering
\caption{V2 AgentState Key Fields}
\label{tab:v2_state}
\begin{tabular}{lll}
\toprule
\textbf{Field} & \textbf{Type} & \textbf{Semantics} \\
\midrule
\texttt{query\_plan}        & QueryPlan               & Reasoning type, sub-queries, difficulty signals \\
\texttt{retrieval\_plans}   & list[RetrievalPlan]     & One plan per sub-query (replaced each loop) \\
\texttt{retrieval\_outputs} & list[RetrievalOutput]   & Retrieved chunks per sub-query (replaced each loop) \\
\texttt{validation\_result} & ValidationResult        & Scores, pass/fail, failure diagnosis \\
\texttt{validated\_evidence}& list[ChunkResult]       & Chunks that passed relevance threshold \\
\texttt{answer}             & AgentAnswer             & Synthesised text with citation IDs \\
\texttt{loop\_count}        & int                     & Re-retrieval iterations completed so far \\
\texttt{trace}              & list[TraceEntry]        & Accumulated event log (append-only) \\
\bottomrule
\end{tabular}
\end{table}


\subsection{Query Analysis Agent}
\label{subsec:v2_qa}

The Query Analysis Agent (QAA) classifies the incoming query into one of six
reasoning types, extracts difficulty signals, and decomposes multi-hop queries
into ordered sub-queries.

\textbf{Classification.} The LLM path sends the raw query to a structured JSON
prompt requesting: \texttt{reasoning\_type} (one of \textit{single\_hop},
\textit{multi\_hop}, \textit{aggregation}, \textit{comparison},
\textit{exception}, \textit{temporal}), a list of \texttt{difficulty\_signals}
(e.g.\ \textit{temporal\_reasoning}, \textit{multi\_document\_dependency},
\textit{exception\_handling}), a boolean \texttt{is\_decomposed}, and a list
of \texttt{sub\_queries} with optional entity slot placeholders.

\textbf{Heuristic fallback.} When no LLM provider is available, the agent
applies a six-pattern regex classifier. The heuristic assigns
\textit{confidence}~$= 0.6$ and produces a single sub-query equal to the raw
query. Regex patterns are ordered by specificity (temporal keywords $>$
exception indicators $>$ multi-hop pronouns $>$ aggregation phrases $>$
comparison phrases $>$ default single\_hop), matching the V1 oracle label
heuristics.

\textbf{Oracle mode.} Evaluation experiments E1--E3 pass
\texttt{use\_oracle=True} with the ground-truth \texttt{reasoning\_type} in
query metadata. The QAA bypasses both LLM and heuristic paths, constructing a
QueryPlan directly from metadata with \textit{confidence}~$= 1.0$. This
isolates the effects of the retrieval loop and scorer from classification
accuracy when measuring RQ8--RQ10.

\textbf{Multi-hop decomposition.} For \textit{multi\_hop} queries, the QAA
produces two or more sub-queries with dependency relationships. Dependent
sub-queries contain entity slot placeholders of the form
\texttt{\{variable\_name\}}, indicating that the value must be extracted from
the results of a prior sub-query. For example, the query ``Which API endpoints
use the same authentication method as the data export service?'' is decomposed
into: (1)~``What authentication method does the data export service use?'' and
(2)~``Which API endpoints use \texttt{\{auth\_method\_q1\}}?'', where the slot
\texttt{\{auth\_method\_q1\}} is to be filled after sub-query~1 is resolved.


\subsection{Planning Agent}
\label{subsec:v2_plan}

The Planning Agent (PA) maps each sub-query to a retrieval strategy using
a three-tier decision process:

\textbf{(1) Re-retrieval override.} If \texttt{loop\_count}~$> 0$ and the
Validation Agent has supplied a \texttt{suggested\_strategy}, that strategy
takes precedence over all other mechanisms. This ensures that failure
diagnosis feedback propagates into the next retrieval attempt.

\textbf{(2) V1 rule planner.} The same seven rules (R01--R07) from V1
(Table~\ref{tab:app_rules}) are applied to the sub-query's reasoning type,
difficulty signals, and document category. Each rule carries a
\textit{confidence} score (0.0--1.0) reflecting the empirical reliability of
its strategy selection. Rules R01--R05 fire with confidence $\geq 0.75$ and
are used directly without LLM consultation.

\textbf{(3) LLM fallback.} Rules R06 (confidence 0.70) and R07 (confidence
0.60) are below the \texttt{LLM\_FALLBACK\_THRESHOLD} of 0.75. When these
rules fire and an LLM provider is available, the agent sends the query plan
to the same structured prompt used in V1 (Appendix~\ref{app:prompt}) and
uses the LLM-selected strategy if it is a valid member of
\{\textit{bm25}, \textit{dense}, \textit{hybrid}\}. If the LLM call fails,
the rule result is used as fallback.


\subsection{Retrieval Agent}
\label{subsec:v2_retrieval}

The Retrieval Agent (RA) executes retrieval plans using the MCP
tool layer \citep{mcp2024} — an in-process server exposing three
retrieval tools (\texttt{bm25\_search}, \texttt{dense\_search},
\texttt{hybrid\_search}) that wrap the pre-built SEKD indexes.

The RA selects an execution mode based on the QueryPlan's reasoning type:

\textbf{Sequential (multi\_hop).} Sub-queries are topologically sorted by
dependency. Each sub-query is executed in order; before executing a dependent
sub-query, entity slots in its query text are filled with values extracted
from the top results of the dependency sub-query. Slot filling uses an LLM
extraction prompt when a provider is available, or a heuristic fallback that
substitutes the first 80 characters of the top-ranked chunk's text.

\textbf{Parallel-merged (aggregation).} All sub-queries are executed
independently. Their chunks are pooled, deduplicated by \texttt{chunk\_id}
(keeping the copy with the highest score), and returned as a single
\texttt{RetrievalOutput} with \texttt{sub\_query\_id}~$=$~``merged''.

\textbf{Independent (all other types).} Each sub-query is executed
independently; results are returned as separate \texttt{RetrievalOutput}
objects for per-sub-query validation scoring.


\subsection{Validation Agent}
\label{subsec:v2_validation}

The Validation Agent (VA) is the core feedback mechanism of the V2 pipeline.
It scores retrieved chunks for relevance, makes a pass/fail decision, and --
on failure -- diagnoses \emph{why} retrieval failed and what should change on
the next attempt.

\subsubsection*{Scorer Implementations}

The VA implements five interchangeable relevance scorers
(Table~\ref{tab:v2_scorers}) selected via the \texttt{scorer} configuration
parameter.

\begin{table}[H]
\centering
\caption{V2 Validation Agent Scorer Implementations}
\label{tab:v2_scorers}
\begin{tabular}{lll}
\toprule
\textbf{Scorer} & \textbf{Signal} & \textbf{Cost} \\
\midrule
\texttt{rank\_proxy}  & Score $= \max(0,\, 1 - (r{-}1) \cdot 0.1)$ for rank $r$ & Zero \\
\texttt{bm25}         & Normalised BM25 score (chunk vs.\ sub-query text) & Negligible \\
\texttt{dense}        & \texttt{chunk.score} (cosine similarity, already in $[0,1]$) & Zero \\
\texttt{adaptive}     & Routes to bm25/dense/avg by \texttt{chunk.strategy\_used} & Negligible \\
\texttt{llm}          & GPT-4o-mini batch scoring (5 chunks/call) & $\sim$\$0.001/10 chunks \\
\bottomrule
\end{tabular}
\end{table}

The default scorer for evaluation experiments E1--E4 is \texttt{adaptive}.
BM25 normalisation uses the formula $s_i = f_i^{(q)} / \max_j f_j^{(q)}$,
where $f_i^{(q)}$ is the raw BM25 score of chunk $i$ for query $q$ and the
maximum is taken over all chunks in the corpus. The dense scorer uses the
cosine similarity stored in \texttt{chunk.score}, which is already bounded
in $[0, 1]$ by the vector index. The adaptive scorer routes each chunk to the
BM25 or dense scorer based on its \texttt{strategy\_used} field, and averages
both scores for hybrid-strategy chunks.

\subsubsection*{Pass Condition and Thresholds}

A chunk is considered \emph{relevant} if its relevance score exceeds the
reasoning-type-specific threshold $\theta_{\tau}$. Validation passes when
the number of relevant chunks $|\mathcal{R}|$ meets the minimum evidence
requirement:

\begin{equation}
\text{pass} \iff |\mathcal{R}| \geq \delta, \quad
|\mathcal{R}| = |\{c \in \mathcal{C} : \text{score}(c, q_c) \geq \theta_{\tau_c}\}|
\label{eq:v2_pass}
\end{equation}

where $\mathcal{C}$ is the set of retrieved chunks, $q_c$ is the sub-query
text associated with chunk $c$, $\tau_c$ is its reasoning type, and $\delta = 3$
is the minimum passed-chunk count. The per-type thresholds reflect the
empirically observed score distributions:

\begin{table}[H]
\centering
\caption{Per-Reasoning-Type Relevance Thresholds}
\label{tab:v2_thresholds}
\begin{tabular}{ll}
\toprule
\textbf{Reasoning Type} & \textbf{Threshold $\theta$} \\
\midrule
single\_hop, comparison, aggregation & 0.30 \\
exception                            & 0.22 \\
multi\_hop                           & 0.20 \\
temporal                             & 0.20 \\
\bottomrule
\end{tabular}
\end{table}

Lower thresholds for \textit{temporal}, \textit{exception}, and
\textit{multi\_hop} queries reflect the sparse-signal or low-overlap nature
of these query types: temporal queries rely on exact date-token matching that
yields low BM25 normalised scores when the corpus maximum is high; exception
queries contain conditional language with limited lexical overlap with
retrieved chunks; multi-hop slots introduce noise via slot-filled reformulations.

\subsubsection*{Failure Diagnosis}

When validation fails, the VA's \texttt{\_diagnose\_failure} function
analyses the score distribution to classify the failure and recommend a
corrective strategy. Table~\ref{tab:v2_failure} lists the six failure
categories and their associated recommendations.

\begin{table}[H]
\centering
\caption{V2 Failure Diagnosis Categories}
\label{tab:v2_failure}
\begin{tabular}{lll}
\toprule
\textbf{Diagnosis} & \textbf{Condition} & \textbf{Recommendation} \\
\midrule
\texttt{no\_chunks}             & No chunks retrieved & Switch to hybrid \\
\texttt{wrong\_strategy}        & All scores $\approx 0$; oracle $\neq$ current & Switch to oracle strategy \\
\texttt{too\_narrow}            & All scores $\approx 0$; already on hybrid & Broaden query \\
\texttt{missing\_entity}        & Multi-hop: unfilled slot \texttt{\{...\}} in query & Re-slot-fill \\
\texttt{insufficient\_coverage} & Some relevant ($0 < |\mathcal{R}| < 3$); current $\neq$ hybrid & Switch to hybrid \\
\texttt{low\_quality}           & Generic low quality; current $\neq$ hybrid & Switch to hybrid \\
\bottomrule
\end{tabular}
\end{table}

The empirical oracle strategy table ($\tau \mapsto s^*$) is identical to the
V1 rule planner oracle labels: $\textit{temporal} \mapsto \textit{bm25}$,
$\textit{exception} \mapsto \textit{dense}$, $\textit{multi\_hop} \mapsto
\textit{hybrid}$, and \textit{hybrid} for all remaining types.


\subsection{Answer Agent}
\label{subsec:v2_answer}

The Answer Agent (AA) synthesises a textual response from the set of
validated evidence chunks (\texttt{validated\_evidence}).

\textbf{LLM synthesis.} The top-$N$ evidence chunks (default $N = 10$,
sorted by rank) are formatted as a numbered list with category labels:
\texttt{[i] (category) text[:400]}. An LLM provider is prompted with a
grounding instruction (``cite sources with inline \texttt{[N]} markers''),
the question, and the evidence block. The response is parsed for citation
markers with the regex \texttt{\textbackslash[(\textbackslash d+)\textbackslash]},
yielding a list of cited chunk IDs sorted by citation index. The answer's
\textit{confidence} is computed as the mean retrieval score of cited chunks;
\textit{evidence\_coverage} is $|\text{cited}| / |\text{validated}|$.

\textbf{Extractive fallback.} When no LLM provider is available, the top-3
evidence chunks are concatenated with \texttt{[1]}, \texttt{[2]},
\texttt{[3]} prefixes as a simple extractive answer. Confidence is the mean
score of the three chunks.

\medskip
The AA's output (\texttt{AgentAnswer}) is the pipeline's final product. For the
retrieval evaluation in Chapter~\ref{chap:v2_results}, answer quality is not
directly measured; retrieval metrics (Recall@10, MRR, nDCG@10) are computed
from the \texttt{retrieval\_outputs} of the final pipeline state, enabling a
direct comparison with V1 baseline results.


\subsection{Re-Retrieval Loop and State Management}
\label{subsec:v2_loop}

LangGraph's conditional edge mechanism implements the re-retrieval loop.
After the Validation Agent node, the function \texttt{route\_after\_validation}
inspects the state and returns one of three routing decisions:

\begin{itemize}
  \item \texttt{"answer"}: validation passed ($|\mathcal{R}| \geq \delta$) ---
        proceed to Answer Agent.
  \item \texttt{"re\_retrieve"}: validation failed and
        \texttt{loop\_count}~$<$~\texttt{MAX\_LOOPS} (default 3) --- route to
        the LoopCounter node (which increments \texttt{loop\_count} and writes
        a trace entry), then back to the Planning Agent.
  \item \texttt{"end"}: loop limit reached, unrecoverable error, or no
        \texttt{validation\_result} --- terminate the graph without answer
        generation.
\end{itemize}

On re-entry into the Planning Agent, the \texttt{validation\_result}
from the previous iteration is visible in the state; the re-retrieval
override mechanism (Section~\ref{subsec:v2_plan}) uses
\texttt{validation\_result.suggested\_strategy} to select a different
retrieval strategy than the one that failed. In this way, each loop
iteration can correct a strategy mismatch identified by the failure
diagnosis system.


\section{Evaluation Framework}
\label{sec:v2_eval_framework}

\subsection{Experimental Configurations}
\label{subsec:v2_experiments}

Four experiments are designed to ablate the independent contributions of
the QA Agent classification accuracy, the Validation Agent scorer, and
the re-retrieval loop. Table~\ref{tab:v2_experiments} summarises the
configurations.

\begin{table}[H]
\centering
\caption{V2 Experiment Configurations}
\label{tab:v2_experiments}
\begin{tabular}{lllll}
\toprule
\textbf{Exp.} & \textbf{Name} & \textbf{QA Mode} & \textbf{Scorer} & \textbf{Max Loops} \\
\midrule
E1 & \texttt{v2\_oracle}           & Oracle           & Adaptive    & 3 \\
E2 & \texttt{v2\_oracle\_noloop}   & Oracle           & Adaptive    & 0 \\
E3 & \texttt{v2\_oracle\_rankproxy}& Oracle           & Rank proxy  & 3 \\
E4 & \texttt{v2\_heuristic}        & Heuristic        & Adaptive    & 3 \\
\bottomrule
\end{tabular}
\end{table}

\textbf{E1 (v2\_oracle)} is the primary V2 configuration: perfect query
classification, adaptive scorer, and up to three re-retrieval loops. This
experiment measures the maximum potential of the agentic pipeline under ideal
classification conditions.

\textbf{E2 (v2\_oracle\_noloop)} disables re-retrieval by setting
\texttt{MAX\_LOOPS}$=0$. Since E2 and E1 share identical QA classification and
retrieval strategy selection, any difference in retrieval metrics is attributable
solely to the re-retrieval loop.

\textbf{E3 (v2\_oracle\_rankproxy)} replaces the adaptive scorer with the
rank-proxy heuristic (score $= 1 - (r-1) \times 0.1$). Since ranks are
determined by the retrieval tool rather than by query-specific relevance,
the rank proxy is insensitive to relevance mismatches. Comparing E3 and E1
isolates the contribution of relevance-based validation.

\textbf{E4 (v2\_heuristic)} uses the regex heuristic classifier (confidence
0.6) in place of oracle classification. All V1 studies used oracle labels;
E4 measures the QA classification error that accumulates in a deployment
scenario where ground-truth reasoning types are unavailable.

\subsection{Research Questions}
\label{subsec:v2_rqs}

The V2 evaluation addresses six additional research questions (RQ8--RQ13),
which are also presented in Section~\ref{sec:v2_rqs_intro}:

\begin{description}
  \item[RQ8] Does the V2 agentic pipeline (E1: oracle QA, adaptive scorer,
             re-retrieval) achieve higher Recall@10 and MRR than the best
             V1 baseline (Rule Planner)?
  \item[RQ9] Does the adaptive re-retrieval loop (E1 vs.\ E2) improve
             retrieval performance when initial validation fails?
  \item[RQ10] Does the adaptive relevance scorer (E1 vs.\ E3) enable more
              accurate re-retrieval trigger decisions than rank-proxy scoring?
  \item[RQ11] Does oracle query classification (E1 vs.\ E4) outperform
              heuristic classification, and by how much?
  \item[RQ12] How does V2 pipeline performance vary across the six reasoning
              types and six document categories?
  \item[RQ13] What are the primary failure modes of the V2 pipeline, and how
              does the failure diagnosis system characterise them?
\end{description}

\subsection{Metrics and Comparison}
\label{subsec:v2_metrics}

Retrieval metrics (Recall@K, MRR, nDCG@K, Hit@K for $K \in \{1,3,5,10\}$)
are computed from the \texttt{retrieval\_outputs} of the final pipeline state
using the same evaluation infrastructure as V1 (Section~\ref{subsec:eval_metrics}).
For experiments with re-retrieval ($\texttt{MAX\_LOOPS} > 0$), the final
state's \texttt{retrieval\_outputs} represents the results of the last
retrieval attempt --- the pipeline's best delivered result set.
Chunks are deduplicated by \texttt{chunk\_id} (keeping the highest score)
and re-ranked by score descending before metric computation.

Pipeline statistics (validation pass rate, loop count distribution, failure
reason distribution) are reported from the per-query \texttt{trace} field,
providing a window into the pipeline's internal decision trajectory that has
no analogue in V1.

All V2 experiments are evaluated on the same 380 answerable queries used in
V1, enabling direct metric comparison. The script
\texttt{scripts/run\_v2\_eval.py} implements all four experiments and writes
per-experiment \texttt{metrics.json}, \texttt{retrieval\_results.jsonl},
\texttt{pipeline\_stats.json}, and a unified \texttt{outputs/evaluation/v2/comparison.json}.
